% !TeX program = xelatex
% !TeX encoding = UTF-8
% 在 out 目录中运行：xelatex -interaction=nonstopmode -halt-on-error ../main.tex

\PassOptionsToPackage{no-math}{fontspec}
\documentclass[
  12pt,
  UTF8,
  a4paper,
  fontset=fandol,
  oneside,
  openany
]{ctexrep}

\usepackage{amsmath}
\usepackage{caption}
\usepackage{circuitikz}
\usepackage[
  bottom=2.54cm,
  left=3.18cm,
  right=3.18cm,
  top=2.54cm
]{geometry}
\usepackage[hidelinks,unicode]{hyperref}
\usepackage{unicode-math}
\usepackage{xcolor}

\definecolor{reviewgray}{gray}{0.42}
\DeclareCaptionFont{wuhao}{\zihao{5}}
\DeclareCaptionLabelFormat{figurelabel}{图#2}
\DeclareCaptionLabelSeparator{figurespace}{\hspace{1em}}
\captionsetup[figure]{
  font=wuhao,
  justification=centering,
  labelformat=figurelabel,
  labelsep=figurespace
}
\hypersetup{pdftitle={电力电子技术笔记}}
\linespread{1.3}
\pagestyle{plain}
\raggedbottom{}
\setcounter{secnumdepth}{5}
\setcounter{tocdepth}{3}
\setlength{\parindent}{2em}
\setlength{\parskip}{0pt}
\setmathfont{XITS Math}
\setmathfont[version=bold]{XITS Math}

\newcommand{\reviewnote}[1]{%
  \par\smallskip
  {\color{reviewgray}\small\noindent【校核建议：#1】\par}
  \smallskip
}
\renewcommand{\thechapter}{\arabic{chapter}}
\renewcommand{\theequation}{\arabic{chapter}-\arabic{equation}}
\renewcommand{\thefigure}{\arabic{chapter}-\arabic{figure}}
\renewcommand{\theparagraph}{\thesubsubsection.\arabic{paragraph}}
\renewcommand{\thesection}{\thechapter.\arabic{section}}
\renewcommand{\thesubparagraph}{\theparagraph.\arabic{subparagraph}}
\renewcommand{\thesubsection}{\thesection.\arabic{subsection}}
\renewcommand{\thesubsubsection}{\thesubsection.\arabic{subsubsection}}

\ctexset{
  chapter={
    afterindent=true,
    aftername=\quad,
    afterskip=1.5ex,
    beforeskip=0pt,
    format=\heiti\centering\zihao{2}\bfseries,
    name={第,章},
    number=\chinese{chapter}
  },
  paragraph={
    afterindent=true,
    aftername=\hspace{0.5em},
    afterskip=0ex,
    beforeskip=0.5ex,
    format=\songti\raggedright\zihao{-4},
    indent=0pt,
    name={,},
    number=\theparagraph,
    runin=false
  },
  section={
    afterindent=true,
    aftername=\hspace{0.5em},
    afterskip=1ex,
    beforeskip=1ex,
    format=\heiti\raggedright\zihao{3}\bfseries,
    name={,},
    number=\thesection
  },
  subparagraph={
    afterindent=true,
    aftername=\hspace{0.5em},
    afterskip=0ex,
    beforeskip=0.5ex,
    format=\songti\raggedright\zihao{-4},
    indent=0pt,
    name={,},
    number=\thesubparagraph,
    runin=false
  },
  subsection={
    afterindent=true,
    aftername=\hspace{0.5em},
    afterskip=0ex,
    beforeskip=0.3cm,
    format=\heiti\raggedright\zihao{-4}\bfseries,
    name={,},
    number=\thesubsection
  },
  subsubsection={
    afterindent=true,
    aftername=\hspace{0.5em},
    afterskip=0ex,
    beforeskip=0.3cm,
    format=\songti\raggedright\zihao{-4}\bfseries,
    name={,},
    number=\thesubsubsection
  }
}

\IfFontExistsTF{SimHei}{
  \setCJKfamilyfont{zhhei}[AutoFakeBold=2]{SimHei}
  \setCJKsansfont[AutoFakeBold=2]{SimHei}
}{}
\IfFontExistsTF{SimSun}{
  \setCJKfamilyfont{zhsong}[AutoFakeBold=2]{SimSun}
  \setCJKmainfont[AutoFakeBold=2]{SimSun}
}{}
\IfFontExistsTF{Times New Roman}{
  \setmainfont{Times New Roman}
}{
  \setmainfont{TeX Gyre Termes}
}{}





% ---------- 正文 ----------
\begin{document}

\chapter{变换器与稳态分析}

\section{概述}

\subsection{结构}

\begin{figure}[!htbp]
  \centering
  \begin{tikzpicture}[
    block/.style={
        align=center,
        draw,
        minimum height=0.9cm,
        minimum width=3.2cm
      },
    flow/.style={-{Latex[length=2.2mm]}, line width=0.55pt},
    every node/.style={font=\small}
    ]
    \node[block] (converter) at (0,1.05) {变换器};
    \node[block] (controller) at (0,-0.75) {控制器};
    \coordinate (input-tap) at (-3.2,1.05);
    \coordinate (output-tap) at (3.2,1.05);
    \draw[flow] (-5.0,1.05) -- node[above]{功率输入} (converter.west);
    \draw[flow] (converter.east) -- node[above]{功率输出} (5.0,1.05);
    \draw[flow] (input-tap) |- (controller.west);
    \draw[flow] (output-tap) |- (controller.east);
    \fill (input-tap) circle[radius=1.4pt];
    \fill (output-tap) circle[radius=1.4pt];
    \node[above] at (-2.35,-0.75) {前馈};
    \node[above] at (2.35,-0.75) {反馈};
    \draw[flow] (controller.north) -- node[right]{控制信号输入} (converter.south);
    \draw[flow] (0,-1.85) -- node[right]{参考量} (controller.south);
  \end{tikzpicture}
  \vspace{-0.2cm}
  \caption{变换器系统的基本结构}
  \vspace{-0.5cm}
  \label{fig:converter-system-structure}
\end{figure}

功率电路变换器类型：脉宽调制（PWM）变换器、准谐振／多谐振变换器、软开关变换器；控制电路类型：模拟电路和数字电路。

\subsection{指标}

效率与损耗为
\begin{equation}
  \eta=\frac{P_{\mathrm{out}}}{P_{\mathrm{in}}},\qquad
  P_{\mathrm{loss}}=P_{\mathrm{in}}-P_{\mathrm{out}}.
\end{equation}

损耗来源：开关器件有通态损耗、开关损耗、驱动损耗；电感、变压器线圈有铜损（线路发热）、铁损（铁芯磁滞、涡流）、磁损（线圈漏磁）；电容损耗，如等效串联电阻、漏电等。

\subsection{器件种类}

阻、感、容、晶。

\begin{figure}[!htbp]
  \centering
  \begin{circuitikz}[scale=1.0, transform shape, every node/.style={font=\small}]
    \ctikzset{
      bipoles/length=0.85cm,
      transistors/scale=1.5,
      tripoles/npn/base width=0.42,
      tripoles/npn/width=0.48,
      tripoles/nigfete/base width=0.32,
      tripoles/nigfete/gate width=0.40
    }
    \begin{scope}[scale around={1.15:(0.7,0.2)}]
      \draw (0,0.2) to[american resistor] (1.4,0.2);
    \end{scope}
    \begin{scope}[scale around={1.15:(0.7,-0.8)}]
      \draw (0,-0.8) to[european resistor] (1.4,-0.8);
    \end{scope}
    \begin{scope}[scale around={1.15:(3.7,0.2)}]
      \draw (3.0,0.2) to[L] (4.4,0.2);
    \end{scope}
    \begin{scope}[scale around={1.15:(3.7,-0.8)}]
      \ctikzset{quadpoles style=inline}
      \node[transformer core, rotate=90] (t) at (3.7,-0.8) {};
      \node[circ] at (t.outer dot A1) {};
      \node[circ] at (t.outer dot B1) {};
    \end{scope}
    \begin{scope}[scale around={1.15:(6.7,-0.3)}]
      \draw (6.0,-0.3) to[C] (7.4,-0.3);
    \end{scope}
    \begin{scope}[scale around={1.15:(9.7,-0.2)}]
      \node[npn, nobase, rotate=90, anchor=center] (bjt) at (9.7,-0.2) {};
      \draw (bjt.nobase) -- ++(0,-0.22);
    \end{scope}
    \begin{scope}[scale around={1.15:(12.9,-0.2)}]
      \node[nigfete, nogate, rotate=90, anchor=center] (mos) at (12.9,-0.2) {};
      \draw (mos.inner down |- mos.nogate) -- ++(0,-0.22);
    \end{scope}
    \node[anchor=north] at (0.7,-1.35) {电阻};
    \node[anchor=north] at (3.7,-1.35) {电感};
    \node[anchor=north] at (6.7,-1.35) {电容};
    \node[anchor=north, align=center] at (9.7,-1.35) {晶体管（线性）};
    \node[anchor=north, align=center] at (12.9,-1.35) {晶体管（开关）};
  \end{circuitikz}
  \vspace{-0.2cm}
  \caption{常见器件}
  \vspace{-0.5cm}
  \label{fig:power-electronic-devices}
\end{figure}

电阻、线性半导体器件不常用于传输电路本身的大功率场景，应尽量减少主功率通道中的损耗性器件，以避免不必要的耗散。

\subsection{理想开关的功率损耗}

理想本身开关不消耗功率：开通／闭合时 $v(t)=0$；关断／打开时 $i(t)=0$。故任意时刻存在
\begin{equation}
  p(t)=v(t)i(t)=0.
  \label{eq:ideal-switch-power-loss}
\end{equation}

\subsection{输出滤波}

低通滤波器（LPF）的应用：LC 输出滤波器保留所需的直流分量或交流基波，并衰减开关频率及其谐波，其电容与电感等储能器件的能量变化可参考“蓄水池”。在功率、纹波和截止频率等指标相同的条件下，开关频率越高，所需的 $L$、$C$ 的值越小，设备总重也更小。

\subsection{其他知识}

\subsubsection{强电与弱电之分}

传统强弱电分类存在不足，电力电子为强电，但也有所谓“强电”的功率不“强”，比如USB插头转换器；通信为弱电，但也有所谓“弱电”的功率不“弱”，比如军用微波天线。再者，还需要考虑电压等级、用途等等，不能简单划分。

\subsubsection{电力电子的应用}

日常应用：含 USB 等的排插、充电宝、无线充电、适配器、电脑电源、电车充电桩；
信息技术设备：PC主板、超算、基站、物联网；
交通：高铁、磁悬浮、船舶、航空航天；
军工。

\subsubsection{电力电子主要负载与消费端}

电机，尤其是交—直—交变频；电脑，采用交—直—直结构，中间含 DC link。

\subsubsection{能源与电网}

传统清洁／非清洁能源的一次电路与电力电子关系较弱，风能、太阳能等新能源与电力电子强相关，涉及并网与调度要求。



\newpage{}

\section{稳态分析原则}

\subsection{分析前提}

本节分析中，只分析理想器件在固定开关频率、周期稳态、连续导通模式（CCM）、满足小纹波近似的条件下的情形，并只针对纯电阻负载进行分析。进入断续导通模式（DCM）或计及损耗后，部分结论需要重推。

\subsection{平均值}

以 Buck 变换器的开关节点电压 $v_{\mathrm{s}}(t)$ 为例，
\begin{equation*}
  \langle v_{\mathrm{s}}\rangle
  =\frac{1}{T_{\mathrm{s}}}\int_0^{T_{\mathrm{s}}}v_{\mathrm{s}}(t)\,\mathrm{d}t
  =\frac{1}{T_{\mathrm{s}}}\bigl(DT_{\mathrm{s}}V_{\mathrm{in}}\bigr)
  =DV_{\mathrm{in}}.
\end{equation*}

\subsection{谐波与低通滤波}

谐波：均值以外的、由傅里叶级数展开得到的正弦分量。低通滤波：滤除 $f_{\mathrm{s}}\ne0$ 的谐波【原文表述】。

\reviewnote{建议改为“低通滤波器主要衰减位于 $nf_{\mathrm{s}}$（$n=1,2,\ldots$）附近的开关谐波”。实际截止频率还需兼顾所需动态带宽，不能简单理解为滤除所有非零频率分量。}

\subsection{DC--DC 变换器的变换比}

一般写为
\begin{equation*}
  V_{\mathrm{o}}=V_{\mathrm{in}}M(D).
\end{equation*}
几种基本变换器为
\begin{equation*}
  \text{Buck: }M(D)=D,\qquad
  \text{Boost: }M(D)=\frac{1}{1-D},\qquad
  \text{Buck--Boost: }M(D)=-\frac{D}{1-D}.
\end{equation*}

\subsection{小纹波近似}

\begin{equation*}
  v_{\mathrm{o}}(t)=V_{\mathrm{o}}+v_{\mathrm{ripple}}(t).
\end{equation*}
当 $\lVert v_{\mathrm{ripple}}\rVert\ll V_{\mathrm{o}}$ 时，有 $v_{\mathrm{o}}(t)\approx V_{\mathrm{o}}$，即忽略纹波。

\reviewnote{可进一步说明 $v_{\mathrm{ripple}}(t)$ 的周期平均值为零，并明确范数或改用峰峰值／幅值与 $V_{\mathrm{o}}$ 的比值表示“小”。}

\section{Buck 变换器的稳态分析}

\subsection{两个开关状态}

【图：Buck 变换器。输入电压源 $V_{\mathrm{in}}$ 经过二选一开关和电感 $L$ 接至输出端，输出端为 $C$ 与 $R$ 并联，输出电压为 $v_{\mathrm{o}}(t)$。】

开关接 1 号端时，电感电压
\begin{equation*}
  v_{\mathrm{L}}=V_{\mathrm{in}}-v_{\mathrm{o}}(t)\approx V_{\mathrm{in}}-V_{\mathrm{o}},
\end{equation*}
由本身性质
\begin{equation*}
  v_{\mathrm{L}}(t)=L\frac{\mathrm{d}i_{\mathrm{L}}(t)}{\mathrm{d}t}
  \quad\Longrightarrow\quad
  \frac{\mathrm{d}i_{\mathrm{L}}(t)}{\mathrm{d}t}
  =\frac{v_{\mathrm{L}}(t)}{L}
  \approx\frac{V_{\mathrm{in}}-V_{\mathrm{o}}}{L}.
\end{equation*}
即电感电流以恒定斜率变化（短时间）。

开关接 2 号端时，电感正端接地，
\begin{equation*}
  v_{\mathrm{L}}=-v_{\mathrm{o}}\approx -V_{\mathrm{o}},
  \qquad
  \frac{\mathrm{d}i_{\mathrm{L}}(t)}{\mathrm{d}t}
  \approx-\frac{V_{\mathrm{o}}}{L}.
\end{equation*}
同上，斜率为负且绝对值不一定相同。

\subsection{稳态波形与电感电流纹波}

【图：一个周期内的 Buck 波形。$v_{\mathrm{L}}(t)$ 在 $DT_{\mathrm{s}}$ 内为 $V_{\mathrm{in}}-V_{\mathrm{o}}$，在 $D'T_{\mathrm{s}}$ 内为 $-V_{\mathrm{o}}$；$i_{\mathrm{L}}(t)$ 围绕平均值 $I_{\mathrm{L}}$ 呈三角波，峰峰值为 $2\Delta i_{\mathrm{L}}$。】

电感电流纹波满足
\begin{align*}
  2\Delta i_{\mathrm{L}}
    & =\dot{i}_{\mathrm{L}}DT_{\mathrm{s}}
  =\frac{V_{\mathrm{in}}-V_{\mathrm{o}}}{L}DT_{\mathrm{s}},               \\
  \Delta i_{\mathrm{L}}
    & =\frac{V_{\mathrm{in}}-V_{\mathrm{o}}}{2L}DT_{\mathrm{s}},          \\
  L & =\frac{V_{\mathrm{in}}-V_{\mathrm{o}}}{2\Delta i_{\mathrm{L}}}DT_{\mathrm{s}}.
\end{align*}
电感值增大，纹波减小。【旁注：其余参数已知时，据此设计电感值；原文字迹不清。】

\reviewnote{本文把 $2\Delta i_{\mathrm{L}}$ 定义为峰峰值，因此 $\Delta i_{\mathrm{L}}$ 是纹波幅值。后续所有 $\Delta i$、$\Delta v$ 都应沿用这一约定；很多教材则直接用 $\Delta i_{\mathrm{L}}$ 表示峰峰值。}

疑问：稳态电路分析仅针对 DC--DC 吗？AC--DC 呢？其他呢？为何一定要用电感？电容不行吗？电压纹波是否也可分析？

\reviewnote{伏秒平衡和电荷平衡并不限于 DC--DC，而适用于含电感、电容且达到周期稳态的开关变换器。AC--DC 和 DC--AC 也能使用，但需选定合适的开关周期、基波周期和状态变量。电容电压纹波同样可以分析。}

\subsection{启动瞬态与稳态界定}

【图：$i_{\mathrm{L}}(t)$ 从零逐周期上升，经历启动瞬态后进入周期性三角波稳态；图中标有上升斜率 $(V_{\mathrm{in}}-v_{\mathrm{o}}(t))/L$ 与下降斜率 $-v_{\mathrm{o}}(t)/L$。】

当
\begin{equation*}
  i_{\mathrm{L}}\bigl((n+1)T_{\mathrm{s}}\bigr)=i_{\mathrm{L}}(nT_{\mathrm{s}})
\end{equation*}
时，认为进入稳态。瞬态需由微分方程【迭代或求解】。

\reviewnote{更完整的周期稳态条件是所有状态变量在相同相位满足 $\boldsymbol{x}(t+T_{\mathrm{s}})=\boldsymbol{x}(t)$；这里只检查 $i_{\mathrm{L}}$ 是针对当前简化示例。}

\subsection{电感伏秒平衡}

由
\begin{equation*}
  v_{\mathrm{L}}(t)=L\frac{\mathrm{d}i_{\mathrm{L}}(t)}{\mathrm{d}t}
\end{equation*}
可得
\begin{align*}
  \int_0^{T_{\mathrm{s}}}\,\mathrm{d}i_{\mathrm{L}}(t)
   & =\frac{1}{L}\int_0^{T_{\mathrm{s}}}v_{\mathrm{L}}(t)\,\mathrm{d}t, \\
  i_{\mathrm{L}}(T_{\mathrm{s}})-i_{\mathrm{L}}(0)
   & =\frac{1}{L}\int_0^{T_{\mathrm{s}}}v_{\mathrm{L}}(t)\,\mathrm{d}t.
\end{align*}
稳态时
\begin{equation*}
  0=\frac{1}{L}\int_0^{T_{\mathrm{s}}}v_{\mathrm{L}}(t)\,\mathrm{d}t,
  \qquad
  \frac{1}{T_{\mathrm{s}}}\int_0^{T_{\mathrm{s}}}v_{\mathrm{L}}(t)\,\mathrm{d}t
  =\langle v_{\mathrm{L}}\rangle=0.
\end{equation*}
设 $A$ 为一个周期内 $v_{\mathrm{L}}(t)$ 与 $t$ 轴所围波形的有符号面积，则
\begin{align*}
  0=A
   & =\int_0^{T_{\mathrm{s}}}v_{\mathrm{L}}(t)\,\mathrm{d}t            \\
   & =(V_{\mathrm{in}}-V_{\mathrm{o}})(DT_{\mathrm{s}})+(-V_{\mathrm{o}})(D'T_{\mathrm{s}}).
\end{align*}
平均电压为
\begin{equation*}
  \langle v_{\mathrm{L}}\rangle
  =\frac{A}{T_{\mathrm{s}}}
  =D(V_{\mathrm{in}}-V_{\mathrm{o}})+D'(-V_{\mathrm{o}}).
\end{equation*}
由 $D'=1-D$ 得
\begin{equation*}
  V_{\mathrm{o}}=DV_{\mathrm{in}}.
\end{equation*}

\subsection{电容电荷平衡}

由
\begin{equation*}
  i_{\mathrm{C}}(t)=C\frac{\mathrm{d}v_{\mathrm{C}}(t)}{\mathrm{d}t}
\end{equation*}
可得
\begin{equation*}
  v_{\mathrm{C}}(T_{\mathrm{s}})-v_{\mathrm{C}}(0)=\frac{1}{C}\int_0^{T_{\mathrm{s}}}i_{\mathrm{C}}(t)\,\mathrm{d}t.
\end{equation*}
稳态时
\begin{equation*}
  0=\frac{1}{T_{\mathrm{s}}}\int_0^{T_{\mathrm{s}}}i_{\mathrm{C}}(t)\,\mathrm{d}t
  =\langle i_{\mathrm{C}}\rangle.
\end{equation*}

对于电感，两端 $v_{\mathrm{L}}$ 为激励，$i_{\mathrm{L}}$ 为响应；对于电容，$i_{\mathrm{C}}$ 为激励，$v_{\mathrm{C}}$ 为响应。要使响应达到周期性稳态，应使相应激励在一个周期内的平均值为零【原文表述】。

\reviewnote{“激励—响应”是本次建模所选的观察方式，并非绝对因果关系。核心结论应写成：周期稳态下电感平均电压为零、电容平均电流为零。}

\section{Boost 变换器}

\subsection{等效电路与状态方程}

【图：Boost 变换器。输入电压源串联电感，开关节点经开关接地，并经二极管接到 $C\parallel R$ 输出端。旁注：二极管导通时，电感末端可连接 $V_{\mathrm{in}}$ 正负极【字迹不清】。】

接状态 1 时，
\begin{equation*}
  v_{\mathrm{L}}(t)=V_{\mathrm{in}},
  \qquad
  i_{\mathrm{C}}(t)=-\frac{v_{\mathrm{o}}(t)}{R}\approx-\frac{V_{\mathrm{o}}}{R}.
\end{equation*}
接状态 2 时，
\begin{equation*}
  v_{\mathrm{L}}(t)=V_{\mathrm{in}}-V_{\mathrm{o}},
  \qquad
  i_{\mathrm{C}}(t)=i_{\mathrm{L}}(t)-\frac{V_{\mathrm{o}}}{R}
  \approx I_{\mathrm{L}}-\frac{V_{\mathrm{o}}}{R}.
\end{equation*}

\subsection{变换比与平均电感电流}

由伏秒平衡，
\begin{align*}
  \frac{1}{T_{\mathrm{s}}}\int_0^{T_{\mathrm{s}}}v_{\mathrm{L}}(t)\,\mathrm{d}t
      & =DV_{\mathrm{in}}+(1-D)(V_{\mathrm{in}}-V_{\mathrm{o}})=0, \\
  V_{\mathrm{o}} & =\frac{1}{1-D}V_{\mathrm{in}}.
\end{align*}
由电荷平衡，
\begin{align*}
  \frac{1}{T_{\mathrm{s}}}\int_0^{T_{\mathrm{s}}}i_{\mathrm{C}}(t)\,\mathrm{d}t
      & =-D\frac{V_{\mathrm{o}}}{R}+(1-D)\left(I_{\mathrm{L}}-\frac{V_{\mathrm{o}}}{R}\right)=0, \\
  V_{\mathrm{o}} & =RI_{\mathrm{L}}(1-D),                                             \\
  I_{\mathrm{L}} & =\frac{V_{\mathrm{o}}}{R(1-D)}
  =\frac{V_{\mathrm{in}}}{R(1-D)^2}.
\end{align*}

\subsection{电感电流与电容电压纹波}

\begin{align*}
  2\Delta i_{\mathrm{L}}
             & =\left.\frac{\mathrm{d}i_{\mathrm{L}}(t)}{\mathrm{d}t}\right|_{0\le t<DT_{\mathrm{s}}}DT_{\mathrm{s}}
  =\frac{V_{\mathrm{in}}}{L}DT_{\mathrm{s}},                                                  \\
  \Delta i_{\mathrm{L}} & =\frac{V_{\mathrm{in}}}{2L}DT_{\mathrm{s}},                                    \\[0.5ex]
  -2\Delta v_{\mathrm{C}}
             & =\left.\frac{\mathrm{d}v_{\mathrm{C}}(t)}{\mathrm{d}t}\right|_{0\le t<DT_{\mathrm{s}}}DT_{\mathrm{s}}
  =-\frac{V_{\mathrm{o}}}{RC}DT_{\mathrm{s}},                                                            \\
  \Delta v_{\mathrm{C}} & =\frac{V_{\mathrm{o}}}{2RC}DT_{\mathrm{s}}.
\end{align*}

【图：Boost 的 $i_{\mathrm{L}}$ 三角波。状态 1 的斜率为 $V_{\mathrm{in}}/L$，状态 2 的斜率为 $(V_{\mathrm{in}}-V_{\mathrm{o}})/L$，峰峰值为 $2\Delta i_{\mathrm{L}}$。】

【图：Boost 的 $v_{\mathrm{C}}$ 纹波。状态 1 的斜率为 $-V_{\mathrm{o}}/(RC)$，状态 2 的斜率为 $I_{\mathrm{L}}/C-V_{\mathrm{o}}/(RC)$，峰峰值为 $2\Delta v_{\mathrm{C}}$。】

疑问：【纹波正负值有影响吗？】默认 $\Delta v_{\mathrm{C}}$、$\Delta i_{\mathrm{L}}$ 为正值，故按增大、减小过程计算【原文大意】。

\section{分析方法总结}

\begin{enumerate}
  \item 画电路图，明确状态数目与等效电路。
  \item 划定时间区间，分【开关过程】，根据器件【通断情况】和电路状态列方程；主要针对 $L$、$C$ 的激励与响应。【括号内原式不清，疑似为稳态平均条件。】
  \item 按实际情况进行小纹波近似，分析伏秒平衡、电荷平衡。
  \item 求电压增益 $M(D)$ 与其他状态变量（$I_{\mathrm{L}}$ 等）。
  \item 画图，求纹波；【由第 2 步各状态中的电压或电流极性判定变化方向，原文字迹不清】。
\end{enumerate}

\section{\texorpdfstring{\'{C}uk}{Cuk} 变换器}

\subsection{电路与两个开关状态}

\reviewnote{本节必须固定参考方向。按原图方向，$V_2$ 和 $I_2$ 可为负值；负号表示实际极性或电流方向与所画参考方向相反，不代表幅值为负。}

【图：\'{C}uk 变换器原电路。$V_{\mathrm{in}}$、$L_1$、$C_1$、$L_2$ 沿上支路相连，中间节点经二选一开关接地，输出端为 $C_2\parallel R$；标有 $v_1$、$v_2$、$i_1$、$i_2$。】

【图：状态 1 等效电路。输入源与 $L_1$ 构成回路；$C_1$、$L_2$ 和输出端构成另一回路。】

【图：状态 2 等效电路。输入源、$L_1$ 和 $C_1$ 构成回路；$L_2$ 与输出端构成另一回路。】

状态 1：
\begin{align*}
  v_{\mathrm{L}1} & =V_{\mathrm{in}},                \\
  v_{\mathrm{L}2} & =-V_1-V_2
  \quad\text{（来源：环路压降 $V_1+v_{\mathrm{L}2}+V_2=0$）}, \\
  i_{\mathrm{C}1} & =i_2,                            \\
  i_{\mathrm{C}2} & =i_2-\frac{V_2}{R}.
\end{align*}

状态 2：
\begin{align*}
  v_{\mathrm{L}1} & =V_{\mathrm{in}}-V_1, \\
  v_{\mathrm{L}2} & =-V_2,                \\
  i_{\mathrm{C}1} & =i_1,                 \\
  i_{\mathrm{C}2} & =i_2-\frac{V_2}{R}.
\end{align*}

以上使用小纹波近似：$v_1$、$v_2$、$i_1$、$i_2$ 以相应平均量代替。

\subsection{平均量关系}

伏秒平衡：
\begin{align}
  \langle v_{\mathrm{L}1}\rangle
   & =DV_{\mathrm{in}}+D'(V_{\mathrm{in}}-V_1)=0,\label{eq:cuk-vl1} \\
  \langle v_{\mathrm{L}2}\rangle
   & =D(-V_1-V_2)+D'(-V_2)=0.\label{eq:cuk-vl2}
\end{align}
电荷平衡：
\begin{align}
  \langle i_{\mathrm{C}1}\rangle
   & =DI_2+D'I_1=0,\label{eq:cuk-ic1} \\
  \langle i_{\mathrm{C}2}\rangle
   & =D\left(I_2-\frac{V_2}{R}\right)
  +D'\left(I_2-\frac{V_2}{R}\right)=0.\label{eq:cuk-ic2}
\end{align}

由式~\eqref{eq:cuk-vl1} 得
\begin{equation*}
  V_1=\frac{V_{\mathrm{in}}}{1-D}
  \qquad\text{（同 Boost）}.
\end{equation*}
由式~\eqref{eq:cuk-vl2} 得
\begin{equation*}
  V_2=-\frac{D}{1-D}V_{\mathrm{in}}
  \qquad\text{（同 Buck--Boost）}.
\end{equation*}
由式~\eqref{eq:cuk-ic2} 得
\begin{equation*}
  I_2=\frac{V_2}{R}
  =-\frac{D}{1-D}\frac{V_{\mathrm{in}}}{R}.
\end{equation*}
由式~\eqref{eq:cuk-ic1} 得
\begin{equation*}
  I_1=\frac{D^2V_{\mathrm{in}}}{(1-D)^2R}
  \quad\text{（沿用原笔记的参考方向）}.
\end{equation*}
因此
\begin{equation*}
  M(D)=-\frac{D}{1-D},
\end{equation*}
即发生极性反转：按图示参考方向，输入上正下负，输出上负下正。

疑问：$D>D'$ 时呢？【原笔记写有“$D\ne D'$ 时，$I_2=V_2/R$，欧姆定律”，上下文不清。】

\subsection{纹波斜率与估算}

状态 1 的斜率：
\begin{align*}
  \frac{\mathrm{d}i_1(t)}{\mathrm{d}t} & =\frac{V_{\mathrm{in}}}{L_1},        &
  \frac{\mathrm{d}i_2(t)}{\mathrm{d}t} & =\frac{-V_1-V_2}{L_2},                 \\
  \frac{\mathrm{d}v_1(t)}{\mathrm{d}t} & =\frac{I_2}{C_1},                    &
  \frac{\mathrm{d}v_2(t)}{\mathrm{d}t} & =\frac{I_2}{C_2}-\frac{V_2}{RC_2}=0.
\end{align*}

状态 2 的斜率：
\begin{align*}
  \frac{\mathrm{d}i_1(t)}{\mathrm{d}t} & =\frac{V_{\mathrm{in}}-V_1}{L_1},    &
  \frac{\mathrm{d}i_2(t)}{\mathrm{d}t} & =-\frac{V_2}{L_1}
  \quad\text{【原式分母为 $L_1$，疑似应核对】},                                              \\
  \frac{\mathrm{d}v_1(t)}{\mathrm{d}t} & =\frac{I_1}{C_1},                    &
  \frac{\mathrm{d}v_2(t)}{\mathrm{d}t} & =\frac{I_2}{C_2}-\frac{V_2}{RC_2}=0.
\end{align*}

\reviewnote{由状态 2 中 $L_2$ 所在支路可知，第二式分母应重点核对为 $L_2$，即 $\mathrm{d}i_2/\mathrm{d}t=-V_2/L_2$。此处暂保留手写原式。}

【波形图略。】疑问：$V_2$ 无纹波（理论上）吗？为什么为 0？

\reviewnote{这里的零来自先把 $i_2(t)$ 替换成平均值 $I_2$。平均电容电流为零只说明一个周期内净电荷变化为零，并不表示瞬时电容电流或输出电压纹波恒为零。求纹波时必须保留 $\widetilde{i}_2(t)$。}

原笔记给出的纹波估算为
\begin{align*}
  \Delta i_1 & =\frac{V_{\mathrm{in}}DT_{\mathrm{s}}}{2L_1}, \\
  \Delta i_2 & =\frac{(V_1+V_2)DT_{\mathrm{s}}}{2L_2}
  =\frac{V_{\mathrm{in}}DT_{\mathrm{s}}}{2L_2},              \\
  \Delta v_1 & =-\frac{I_2DT_{\mathrm{s}}}{2C_1}
  =\frac{V_{\mathrm{in}}D^2T_{\mathrm{s}}}{2(1-D)RC_1}.
\end{align*}

\section{输出电压纹波与滤波环节}

【关于 \textup{\'{C}uk} 输出侧电压的一段说明字迹不清：大意为在先前的零阶小纹波模型中，输出侧电压纹波被忽略，且两个开关状态给出的电容平均电流相同，不能直接由该近似得到纹波。】

以 Buck 为例，求电容电流时不能忽略电感电流纹波：
\begin{equation*}
  i_{\mathrm{C}}(t)=\widetilde{i}_{\mathrm{L}}(t)-\widetilde{i}_{\mathrm{o}}(t)
  \approx\widetilde{i}_{\mathrm{L}}(t)-I_{\mathrm{o}}
  =i_{\mathrm{L}}(t)-I_{\mathrm{L}}.
\end{equation*}

\reviewnote{符号建议统一为 $i_{\mathrm{C}}(t)=i_{\mathrm{L}}(t)-i_{\mathrm{o}}(t)$；若负载电流纹波可忽略，则 $i_{\mathrm{o}}(t)\approx I_{\mathrm{o}}=I_{\mathrm{L}}$，从而 $i_{\mathrm{C}}(t)\approx i_{\mathrm{L}}(t)-I_{\mathrm{L}}=\widetilde{i}_{\mathrm{L}}(t)$。原式混用了瞬时量和纹波量。}

【图：一个周期内的电感电流三角纹波；以平均值为基准，半个周期内的三角形阴影面积标为电荷 $Q$。】

由三角形面积与电容电荷变化量，
\begin{equation*}
  \frac{1}{2}\cdot\frac{T_{\mathrm{s}}}{2}\cdot\Delta i_{\mathrm{L}}
  =Q=C(2\Delta v_{\mathrm{o}}),
\end{equation*}
得到
\begin{equation*}
  \Delta v_{\mathrm{o}}=\frac{\Delta i_{\mathrm{L}}T_{\mathrm{s}}}{8C}.
\end{equation*}

\reviewnote{该式沿用“$\Delta$ 为幅值、$2\Delta$ 为峰峰值”的约定，并忽略电容 ESR。实际输出纹波还应叠加近似为 $\Delta i_{\mathrm{C},\mathrm{pp}}\,\mathrm{ESR}$ 的阶跃分量。}

【图：输出电压纹波波形，峰峰值为 $2\Delta v_{\mathrm{o}}$。】

\'{C}uk 变换器略复杂【原文：由于电流中“掺导”，字迹不清】，原笔记写为
\begin{equation*}
  \Delta v_2=\frac{\Delta i_2T_{\mathrm{s}}}{8C}=\cdots
  \quad\text{【原式中的电容下标缺失，待核对】}.
\end{equation*}

\reviewnote{若这里对应输出电容，分母应明确写为 $8C_2$；同时仍需说明小纹波和忽略 ESR 的前提。}

\section{小结}

【原笔记此处仅写有“小结（无）”。】

\end{document}
